\documentclass[11pt,a4paper,twoside]{tesis}
% SI NO PENSAS IMPRIMIRLO EN FORMATO LIBRO PODES USAR
%\documentclass[11pt,a4paper]{tesis}

\usepackage{graphicx}
\usepackage[utf8]{inputenc}
\usepackage[spanish]{babel}
\usepackage[left=3cm,right=3cm,bottom=3.5cm,top=3.5cm]{geometry}
\usepackage[version=3]{mhchem}	%fórmulas químicas
\usepackage{siunitx}			%unidades

\newcommand{\h}{\ce{H^+}}
\newcommand{\oh}{\ce{OH^-}}
\newcommand{\na}{\ce{Na^+}}
\newcommand{\cl}{\ce{Cl^-}}
\newcommand{\kvm}{$\si{\kilo\volt\per\metre}$}
\newcommand{\usec}{$\si{\micro\second}$}

\begin{document}

\graphicspath{graficos}

%%%% CARATULA
\def\titulo{Licenciado }

\def\autor{Mauricio Alfonso}
\def\tituloTesis{Estudio de los Mecanismos Básicos de Electroporación a Través de la Modelación Numérica}
\def\runtitulo{Estudio de los Mecanismos Básicos de Electroporación a Través de la Modelación Numérica}
\def\runtitle{Study of the Basic Mechanisms of Electroporation Through Numeric Modelling}
\def\director{Alejandro Soba}
\def\codirector{Guillermo Marshall}
\def\lugar{Buenos Aires, 2014}
\newcommand{\HRule}{\rule{\linewidth}{0.2mm}}
%
\thispagestyle{empty}

\begin{center}\leavevmode

\vspace{-2cm}

\begin{tabular}{l}
\includegraphics[width=2.6cm]{logofcen.pdf}
\end{tabular}


{\large \sc Universidad de Buenos Aires

Facultad de Ciencias Exactas y Naturales

Departamento de Computaci\'on}

\vspace{6.0cm}

%\vspace{3.0cm}
%{
%\Large \color{red}
%\begin{tabular}{|p{2cm}cp{2cm}|}
%\hline
%& Pre-Final Version: \today &\\
%\hline
%\end{tabular}
%}
%\vspace{2.5cm}

{\huge\bf \tituloTesis}

\vspace{2cm}

{\large Tesis presentada para optar al t\'{\i}tulo de\\
\titulo en Ciencias de la Computaci\'on}

\vspace{2cm}

{\Large \autor}

\end{center}

\vfill

{\large

{Director: \director}

\vspace{.2cm}

{Codirector: \codirector}

\vspace{.2cm}

\lugar
}

\newpage\thispagestyle{empty}


%%%% ABSTRACTS, AGRADECIMIENTOS Y DEDICATORIA
\frontmatter
\pagestyle{empty}
%\begin{center}
%\large \bf \runtitulo
%\end{center}
%\vspace{1cm}
\chapter*{\runtitulo}

\noindent 

Esto hay que reescribirlo todo! La electroporación reversible es un método consistente en la aplicación de pulsos eléctricos de alta intensidad a una célula con el objetivo de permeabilizar su membrana creando poros, y así permitir el ingreso de drogas o moléculas de ADN a su interior. Esto permite tratar tumores con menores cantidades de drogas, reduciendo los efectos secundarios. En este trabajo se simula una célula esférica a la que se le aplica un pulso eléctrico de 20\si{\milli\second} de duración a través de dos electrodos, y se estudia el ingreso al interior de la célula de 4 especies iónicas: el ión hidrógeno (\h), el hidróxido (\oh), el catión sodio (\na) y el cloruro (\cl). Para eso se tiene en cuenta el campo eléctrico producido por los electrodos, la generación y evolución de poros en la membrana celular producto de la diferencia de potencial entre el interior y exterior de la célula , y la migración de las especies mencionadas, producto de la diferencia de potencial.
Las simulaciones se realizaron con el método de elementos finitos sobre mallas bidimensionales que representan el dominio sobre un sistema de coordenadas cilíndricas usando elementos cuadrilaterales.

\bigskip

\noindent\textbf{Palabras claves:} Guerra, Rebelión, Wookie, Jedi, Fuerza, Imperio (no menos de 5).

\cleardoublepage
%\begin{center}
%\large \bf \runtitle
%\end{center}
%\vspace{1cm}
\chapter*{\runtitle}

\noindent 
%La electroporación consiste en la aplicación de pulsos eléctricos de alta intensidad y corta duración con el objetivo de crear poros en la membrana celular, logrando así un aumento de la permeabilización que permite el ingreso de drogas o iones a su interior. 
Electroporation involves the application of electric pulses of high intensity and short duration in order to create pores in the cell membrane, thus achieving increased permeabilization that allows the entry of drugs or ions.
%La utilización de la electroporación en combinación con drogas antitumorales ha demostrado tener una significativa mayor eficacia que la terapia quimioterapéutica convencional, de allí la relevancia de estudios básicos de la interacción campos eléctricos-célula. 
The use of electroporation in combination with antitumor drugs has been shown to have a higher efficiency than conventional chemotherapeutic therapy, hence the relevance of basic studies of the electric field-cell interaction.
%En esta tesis se presenta un nuevo modelo numérico que describe la respuesta eléctrica de la célula, en particular la membrana celular y el transporte iónico a través de la misma, a la aplicación de pulsos eléctricos. 
In this thesis a new numerical model is presented, describing the electrical response of the cell, particularly the cell membrane and ion transport through it, to the application of electrical pulses.
%Se asume una célula esférica sometida a pulsos eléctricos por medio de dos electrodos, constituida por cuatro especies iónicas: el ión hidrógeno (\h), el hidróxido (\oh), el catión sodio (\na) y el cloruro (\cl). 
A spherical cell is assumed, subjected to electrical pulses by means of two electrodes and consisting of four ionic species: hydrogen ion (\h), hydroxide (\oh), sodium cation (\na) and chloride (\cl).
%Para resolver las ecuaciones diferenciales que describen el potencial electrostático y el transporte iónico se usó el método de los elementos finitos en dos dimensiones espaciales en coordenadas cilíndricas. 
To solve the differential equations describing the electrostatic potential and ion transport, the finite element method in two spatial dimensions with cylindrical coordinates is used.
%Las ecuaciones diferenciales que describen la evolución de la población de poros se resuelven por diferencias finitas utilizando el método de Euler. 
The differential equations describing the evolution of the population of pores are solved by finite differences using Euler's method.
%Se utilizó programación distribuida basada en OpenMP para aprovechar al máximo los procesadores multithreading actuales. 
Distributed programming based on OpenMP was used to maximize the usage of current multithreading processors.
%El nuevo modelo teórico introducido permite por primera vez predecir realísticamente la respuesta eléctrica de la célula, en particular el campo eléctrico transmembranal y el transporte iónico (uptake), lo que se evidencia por la excelente correlación entre predicción y mediciones.
The new theoretical model presented realistically predicts for the first time the electrical response of the cell, particularly the electric field and ion transport (uptake), as evidenced by the excellent correlation between predictions and measurements. 

\bigskip

\noindent\textbf{Keywords:} cell membrane, electroporation, transport, finite elements

\cleardoublepage
\tableofcontents

\mainmatter
\pagestyle{headings}

%%%% ACA VA EL CONTENIDO DE LA TESIS

% Si tarda mucho en compilar se puede poner cada capítulo en archivos separados.

% CAP 1 introducción. problemática e introducción del modelo
% CAP 2 descripción del modelo. qué se resuelve, método numérico. Explicación de FEM
% CAP 3 ITV. como se resuelve, resultados (solos, sin poros ni transporte)
% CAP 4 Poros. ecuaciones, como se resuelve, resultados
% CAP 5 Transporte. idem (solo transporte, sin poros)
% CAP 6 Todo acoplado. Resultados. poner snapshots de valores 1-9
% CAP 7 conclusiones

\chapter{Introducción}
% CAP 1 introducción. problemática e introducción del modelo

\chapter{Descripción del Modelo}
% CAP 2 descripción del modelo. qué se resuelve, método numérico. Explicación de FEM
% podría ir teoría de gradientes conjugados, etc
% dividir en sections!

% explicación FEM 







\chapter{Potencial Eléctrico}
% CAP 3 ITV. como se resuelve, resultados (solos, sin poros ni transporte)

\chapter{Generación y Evolución de Poros}
% CAP 4 Poros. ecuaciones, como se resuelve, resultados

\chapter{Transporte de Especies}
% CAP 5 Transporte. idem (solo transporte, sin poros)

\chapter{Modelo Acoplado}
% CAP 6 Todo acoplado. Resultados. poner snapshots de valores 1-9

\chapter{Conclusiones}
% CAP 7 conclusiones

%%%% BIBLIOGRAFIA
\backmatter
%\bibliography{tesis}

\end{document}
